\documentclass[../../../../main.tex]{subfiles}
\begin{document}

\begin{flushright}
\footnotesize\textit{【自動文字起こし・要確認】}
\end{flushright}

\noindent
[\textbf{解}] (イ) 帰納法で示す. $n=2$ の時 成立. $n=k$ での成立を仮定すると
\begin{align*}
x_{k+1} &= \frac{x_k^2 + 2}{3} > \frac{1+2}{3} = 1 \\
x_{k+1} &< \frac{a^2 + 2}{3} < a \quad (\because a^2 - 3a + 2 = (a-2)(a-1) < 0)
\end{align*}
より $n=k+1$ でも成立. 以上から示された.

\bigskip

\noindent
(ロ)
\begin{align*}
|x_{n+1} - 1| &= \frac{1}{3}(x_n^2 + 2) - 1 = \frac{1}{3}(x_n + 1)(x_n - 1) \\
&< \frac{1}{3}(a + 1)(x_n - 1) \quad (\because (\text{イ}))
\end{align*}

\bigskip

\noindent
(ハ) くり返し用いて
\[
0 < x_n - 1 \le \left(\frac{a+1}{3}\right)^{n-1}(a - 1)
\]
よつて はさみうちから $\displaystyle x_n \to 1 \ (n \to \infty)$ \hfill $\mathbin{/\mkern-5mu/}$

\end{document}
